\begin{result}{E.2.28}{The completion preserved by all affine integral homology maps}{Proved}
Let \(F=\langle x_1,\ldots,x_r\rangle\) be a specified discrete
free group of finite rank, and put \(G=\widehat F\). For each
open normal subgroup \(N\subset G\), set
\[
 \begin{aligned}
 K_N&=F\cap N,& A_N&=K_N/[K_N,K_N],\\
 D_N&=\overline{[K_N,K_N]}^{\,G},&
 \Gamma_N&=F/[K_N,K_N],\\
 P_N&=G/D_N,& Q_N&=G/N=F/K_N.
 \end{aligned}
 \tag{E.2.28a}
\]
The group \(\Gamma_N\) embeds in \(P_N\). Define the actual
subgroup and its displayed inverse-limit realization by
\[
 \mathcal H_{\mathrm{int}}
   =\bigcap_N\pi_N^{-1}(\Gamma_N)\subset G,
 \qquad
 \mathcal H_{\mathrm{int}}\simeq\varprojlim_N\Gamma_N,
 \quad \pi_N:G\longrightarrow P_N.
 \tag{E.2.28b}
\]
For an open subgroup \(U\subset G\), write
\(\Lambda_U=(F\cap U)^{\mathrm{ab}}\subset U^{\mathrm{ab}}\).
This is the integral homology lattice of its actual finite cover
of the bouquet of \(r\) circles. For a continuous automorphism
\(\alpha\) of \(G\), the following conditions are equivalent:
\begin{enumerate}
\item The actual induced map \(U^{\mathrm{ab}}\to
\alpha(U)^{\mathrm{ab}}\) carries \(\Lambda_U\) onto
\(\Lambda_{\alpha(U)}\) for every open \(U\).
\item The same assertion holds on the integral lattices \(A_N\)
for one characteristic cofinal chain of open normal subgroups.
\item There is one element \(h\in G\) such that
\((\operatorname{ad}(h)\alpha)(\mathcal H_{\mathrm{int}})
=\mathcal H_{\mathrm{int}}\).
\end{enumerate}
Furthermore, an actual continuous automorphism \(\beta\) of
\(G\) preserves \(\mathcal H_{\mathrm{int}}\) if and only if
\(\beta(x_j)\in\mathcal H_{\mathrm{int}}\) for
\(1\leq j\leq r\).

For the punctured curve of E.2.22 this characterizes preservation
of all \emph{affine} integral homology lattices, up to one actual
inner adjustment. Passing from the proper-homology condition in
E.2.22 to this affine condition requires the vanishing of the
boundary-extension classes identified below.
The discrete-representative and Galois requirements remain open.
\end{result}
\begin{proof}
We prove the embedded quotients, the local conjugator calculation,
and the surjectivity that makes all local conjugators compatible.

\paragraph{1. The actual quotient and completion maps.}
There is an exact discrete extension
\[
 1\longrightarrow A_N\longrightarrow\Gamma_N
   \longrightarrow Q_N\longrightarrow1.
 \tag{E.2.28c}
\]
The finite-index free group \(K_N\) has a finite-rank free
abelianization. The group \(\Gamma_N\) is residually finite:
an element outside \(A_N\) is separated by \(Q_N\); a nonzero
element of \(A_N\) is separated by
\(\Gamma_N/mA_N\) for a suitable positive integer \(m\).
Here \(mA_N\) is normal in \(\Gamma_N\), and the quotient
is finite. Hence
\[
 F\cap D_N=[K_N,K_N],\qquad
 \Gamma_N\hookrightarrow P_N=\widehat{\Gamma_N}.
 \tag{E.2.28d}
\]
The completion equality follows from the finite quotient
universal property: a finite quotient of \(F\) kills
\([K_N,K_N]\) exactly when its extension to \(G\) kills
\(D_N\). The subgroups \(mA_N\) also show that the full
profinite topology of \(A_N\) is induced from \(\Gamma_N\).
Thus the actual completed extension is
\[
 1\longrightarrow\widehat A_N\longrightarrow P_N
       \longrightarrow Q_N\longrightarrow1,
 \qquad N/D_N=\widehat A_N,\qquad
 \Gamma_N\cap\widehat A_N=A_N.
 \tag{E.2.28e}
\]
Since \(K_N\) is dense in \(N\), continuity of the commutator
map identifies \(D_N\) with the closed commutator subgroup
of \(N\).

For \(M\subset N\), the inclusions of commutator subgroups
give surjective maps \(P_M\to P_N\) and
\(\Gamma_M\to\Gamma_N\). Compatible coordinates in the
\(P_N\) define nonempty compatible compact cosets of \(D_N\)
in \(G\). Their intersection is one point, since
\(D_N\subset N\) and the open normal subgroups have trivial
intersection. It follows that
\[
 G\simeq\varprojlim_NP_N,\qquad
 \mathcal H_{\mathrm{int}}\simeq\varprojlim_N\Gamma_N.
 \tag{E.2.28f}
\]
Giving each \(\Gamma_N\) the discrete topology makes the
second limit the completion of \(F\) for the neighborhoods
\([K_N,K_N]\); its projections are onto, because \(F\)
already projects onto every \(\Gamma_N\). This topology
is specified separately from the topology induced from \(G\).
For rank two, E.2.27 supplies an element of
\(\mathcal H_{\mathrm{int}}\) outside every profinite
conjugate of \(F\).

One explicit characteristic cofinal chain is obtained by
intersecting the kernels of all continuous homomorphisms
\(G\to S_a\) for \(a\leq n\). There are finitely many such
homomorphisms for each \(n\), since the images of the specified
generators determine them. The resulting \(N_n\) is open and
characteristic. Regular permutation representations show
cofinality. Every automorphism of \(G\) therefore preserves
\(N_n\) and its closed commutator \(D_{N_n}\). We use this
chain for the construction below; it defines the same subgroup
in (E.2.28b).

\paragraph{2. The image in every original cover lattice.}
For every open \(U\subset G\), including nonnormal ones,
\[
 \operatorname{ab}_U(\mathcal H_{\mathrm{int}}\cap U)
       =\Lambda_U.
 \tag{E.2.28g}
\]
To prove the nontrivial inclusion, take
\(z\in\mathcal H_{\mathrm{int}}\cap U\), choose normal
\(N\subset\operatorname{core}_G(U)\), and choose \(f\in F\)
with \(zf^{-1}\in D_N\). The latter subgroup is contained in
\(U\) and in the kernel of \(\operatorname{ab}_U\).
Thus \(f\in F\cap U\) and
\(\operatorname{ab}_U(z)=\operatorname{ab}_U(f)\).
The reverse inclusion follows from
\(F\cap U\subset\mathcal H_{\mathrm{int}}\cap U\).
As in E.2.27, the full finite-index subgroup topology gives
\(U^{\mathrm{ab}}=\Lambda_U\otimes\widehat\Z\).

An arbitrary inner automorphism also carries these actual
lattices to the corresponding actual lattices. Given \(g\in G\),
choose \(f\in F\cap gU\); then \(g=fu\) with \(u\in U\).
Hence \(gUg^{-1}=fUf^{-1}\), and conjugation by \(u\)
acts trivially on \(U^{\mathrm{ab}}\). The induced homology
map is therefore the discrete change-of-basepoint map
given by \(f\).

\paragraph{3. The exact local conjugator.}
Fix \(N\) in the characteristic chain, and suppose \(\alpha\)
carries \(A_N\) onto itself in \(\widehat A_N=N/D_N\).
Let \(\alpha_N\) be its induced automorphism of \(P_N\).
Both \(\Gamma_N\) and \(\alpha_N(\Gamma_N)\) contain
\(A_N\), intersect \(\widehat A_N\) exactly in \(A_N\),
and map onto \(Q_N\).

The abstract abelian group
\(V_N=\widehat A_N/A_N\) is uniquely divisible.
Multiplication by a positive integer \(m\) is surjective
because a profinite vector has an integral representative
modulo \(m\). It is injective because an integral vector
divisible by \(m\) in every \(\Z_p\) is divisible by \(m\)
in \(\Z\). Thus \(V_N\) is a \(\Q\)-vector space.
The abstract extension
\[
 1\longrightarrow V_N\longrightarrow P_N/A_N
        \longrightarrow Q_N\longrightarrow1
 \tag{E.2.28h}
\]
is split by \(\Gamma_N/A_N\), and
\(\alpha_N(\Gamma_N)/A_N\) is a second complement.
Their difference is a cocycle \(c:Q_N\to V_N\).
The literal averaging formula is
\[
 v=\frac1{|Q_N|}\sum_{q\in Q_N}c(q),\qquad
 c(q)=v-qv.
 \tag{E.2.28i}
\]
Indeed, summing \(c(qs)=c(q)+qc(s)\) over \(s\)
gives the second identity. Conjugation by a lift of
\(-v\) in \(\widehat A_N\) carries the second complement
to the first. Before quotienting, both subgroups contain
\(A_N\), so it carries
\(\alpha_N(\Gamma_N)\) exactly onto \(\Gamma_N\).

Consequently the solution set
\[
 S_N=\{u\in P_N:
   \operatorname{ad}(u)\alpha_N(\Gamma_N)=\Gamma_N\}
 \tag{E.2.28j}
\]
is a nonempty orbit for left multiplication by
\(B_N=\operatorname{Norm}_{P_N}(\Gamma_N)\).
We have used \(P_N/A_N\) only as an abstract group;
no Hausdorff or compactness assertion for this quotient
or for \(S_N\) is needed.

\paragraph{4. The normalizer and its integral equations.}
Put \(Z_N=(\widehat A_N)^{Q_N}\). It is central in \(P_N\).
The exact normalizer is
\[
 B_N=\Gamma_N Z_N.
 \tag{E.2.28k}
\]
To prove this, write an arbitrary \(b\in P_N\) as
\(b=\gamma a\), where \(\gamma\in\Gamma_N\) has the same
finite quotient and \(a\in\widehat A_N\). It normalizes
\(\Gamma_N\) precisely when
\[
 (q-1)a\in A_N\quad(q\in Q_N).
 \tag{E.2.28l}
\]
In an integral basis of \(A_N\), stack the integer matrices
\(q-1\) into one matrix \(M\). The vector \(Ma\) is integral.
An integer linear system \(Mz=c\) with integral \(c\) has
an integral solution whenever it has a solution over
\(\widehat\Z\): Smith normal form reduces this to divisibility
of an integer by each nonzero diagonal entry at every prime,
and to vanishing in the zero rows. Both imply the required
integer conditions. Choose such \(a_0\in A_N\) with
\(Ma_0=Ma\). Then \(a-a_0\in Z_N\), proving (E.2.28k).
The reverse inclusion follows from centrality of \(Z_N\).

\paragraph{5. Surjectivity of the actual transition maps.}
The regular cover of the bouquet of \(r\) oriented circles
has vertices \(Q_N\) and edges \((q,j)\), directed from
\(q\) to \(qx_j\). Its integral cycle module is \(A_N\).
There are no two-cells. An invariant edge chain has one
constant coefficient for each label, and every such chain
has zero boundary. Thus, with actual cycles
\(e_j^N=\sum_{q\in Q_N}(q,j)\),
\[
 A_N^{Q_N}=\bigoplus_{j=1}^r\Z e_j^N,\qquad
 Z_N=\bigoplus_{j=1}^r\widehat\Z e_j^N.
 \tag{E.2.28m}
\]
The completed equality follows from the same permutation
coordinates on the edge module. These are the actual
relation-module maps for the cover graph; compare
\cite[pp.~330--331]{GenBuxErshovRapinchuk2011}.

For \(M\subset N\), the cover map sends each edge to the edge
with its quotient vertex and the same label. Consequently
\[
 e_j^M\longmapsto [N:M]e_j^N.
 \tag{E.2.28n}
\]
Write \(d=[N:M]\). The image of \(Z_M\to Z_N\) is
\(dZ_N\), while \(\Gamma_M\to\Gamma_N\) is onto.
Since \(\Gamma_N\) contains the integral invariant cycles,
these two facts give
\[
 B_M\longrightarrow B_N\quad\text{surjective}.
 \tag{E.2.28o}
\]
Explicitly, for \(\gamma z\in\Gamma_N Z_N\), lift
\(\gamma\) through \(\Gamma_M\). Choose an integral
invariant cycle \(z_0\) representing the coordinates of
\(z\) modulo \(d\), and lift it through \(\Gamma_M\).
Finally lift \(z-z_0\) through \(Z_M\) by (E.2.28n).
The product of these three lifts maps to \(\gamma z\).
Centrality of the invariant part makes this the stated
group-map identity.

\paragraph{6. One compatible conjugator for the whole tower.}
Assume condition 2. The map \(P_M\to P_N\) carries
\(S_M\) into \(S_N\), since both \(\Gamma_M\) and its
\(\alpha_M\)-image project onto their indicated groups
at level \(N\). It is surjective on solution sets.
Given \(s_N\in S_N\), take any \(t_M\in S_M\) and let
\(t_N\) be its projection. Then
\(s_Nt_N^{-1}\in B_N\). Lift this element to
\(b_M\in B_M\) using (E.2.28o). The element
\(b_Mt_M\in S_M\) projects exactly to \(s_N\).

Successive choices along the countable characteristic chain
therefore give compatible coordinates in every \(S_N\).
Equation (E.2.28f) produces one actual \(h\in G\) with
those coordinates. Its inner adjustment \(\beta\) satisfies
\(\beta_N(\Gamma_N)=\Gamma_N\) at every level, hence
\(\beta(\mathcal H_{\mathrm{int}})=\mathcal H_{\mathrm{int}}\).
This proves condition 3. The construction uses surjective
transition maps of the solution sets; it does not apply
compactness to nonclosed sets of conjugators. Compactness
is used only for the compatible closed cosets of \(D_N\)
that define an actual element of \(G\).

Conversely, condition 3 and (E.2.28g) show that \(\beta\)
carries every actual integral cover lattice onto the
corresponding lattice. Removing the inner adjustment
preserves that assertion by paragraph 2. Thus condition 1
follows. Its implication to condition 2 is immediate.
In particular the proved group identity is
\[
 \mathcal I_{\mathrm{aff}}
   =\operatorname{Inn}(G)\,
       \operatorname{Stab}_{\Aut_{\mathrm{cont}}(G)}
                    (\mathcal H_{\mathrm{int}}).
 \tag{E.2.28p}
\]
Here \(\mathcal I_{\mathrm{aff}}\) means exactly condition 1;
the stabilizer is explicitly taken inside actual continuous
automorphisms of \(G\).

\paragraph{7. The finite-generator characterization.}
One implication is immediate since \(F\subset
\mathcal H_{\mathrm{int}}\). Conversely, suppose
\(\beta(x_j)\in\mathcal H_{\mathrm{int}}\) for each
generator. For each characteristic \(N\), this gives
\(\beta_N(\Gamma_N)\subset\Gamma_N\). By (E.2.28e),
\(\beta_N\) carries \(A_N\) into \(A_N\). Its integral
matrix has invertible completion, because the original
\(\beta\) is an automorphism of \(N/D_N\). Its determinant
is consequently an integer unit in every \(\Z_p\), and
equals \(1\) or \(-1\). Thus it maps \(A_N\) onto itself.
Its map on the finite \(Q_N\) is also an automorphism.
A subgroup of \(\Gamma_N\) containing \(A_N\) and mapping
onto \(Q_N\) is all of \(\Gamma_N\). Hence
\(\beta_N(\Gamma_N)=\Gamma_N\) for every \(N\), proving
the assertion by (E.2.28f).
This uses an already given automorphism of \(G\); it does
not assert that arbitrary tuples in
\(\mathcal H_{\mathrm{int}}\) generate \(G\).

\paragraph{8. The exact remaining proper-homology block.}
For a finite affine cover \(Y\) with compactification \(C\)
and cusp set \(D\), the actual integral sequence is
\[
 0\longrightarrow B_Y\longrightarrow H_1(Y,\Z)
     \longrightarrow H_1(C,\Z)\longrightarrow0,\qquad
 B_Y=\Z[D]\Big/\Z\Bigl(\sum_{d\in D}d\Bigr).
 \tag{E.2.28q}
\]
The boundary basis maps to the oriented cusp loops, with
the single sum relation. These are free abelian groups,
so the sequence admits integral splittings. Under two
chosen integral splittings, an actual profinite homology
map preserving the completed boundary submodule has matrix
\[
 \begin{pmatrix}P&T\\0&S\end{pmatrix},\qquad
 T\in\operatorname{Hom}_{\Z}(H_1(C,\Z),B_{Y'})
                          \otimes\widehat\Z.
 \tag{E.2.28r}
\]
When \(P\) and \(S\) are integral isomorphisms, the full
map is integral exactly when \(T\) is integral. Changing
either splitting changes \(T\) by an integral homomorphism,
so its vanishing class modulo the integral Hom group
is independent of the splittings.

For the actual \(\sigma\in\Gamma_{F_{\mathrm{cm}}}\) of
E.2.22, the cyclotomic character is one. Preservation of
cusp inertia gives the integral cusp permutation \(P\),
and E.2.22 prescribes integrality of \(S\). The still
required map is the vanishing of the resulting actual
boundary-extension classes, simultaneously for that one
Galois automorphism over the whole cover tower. Neither
block equation by itself gives this vanishing.

After affine integrality, (E.2.28p) identifies the actual
larger completion that is preserved up to one inner
conjugator. The further discrete representative is still
needed for the specific Galois and peripheral action.
For example, the exact hypotheses of the opened virtual
homology tests in \cite[Theorem~1.2, p.~3; Theorem~1.5,
p.~4]{GenYuksel2025} start, respectively, with an
endomorphism of \(F\) and an element of \(\operatorname{Out}(F)\).
The maps proved here land in the displayed stabilizer
inside \(\Aut(G)\); no representative carrying \(F\)
to \(F\) has been constructed. This states the missing
input to those particular theorems without asserting
that no other mathematical route can supply it.
\end{proof}
